% Generated draft: requires mathematical and source review.
\documentclass[11pt]{article}
\usepackage[margin=25mm]{geometry}
\usepackage{fontspec}
\setmainfont{Latin Modern Roman}
\usepackage{amsmath,amssymb,mathtools,unicode-math}
\setmathfont{Latin Modern Math}
\usepackage{enumitem,needspace,xurl,hyperref}
\setlength{\parindent}{0pt}
\setlength{\parskip}{5pt}
\setlength{\emergencystretch}{3em}
\newcommand{\N}{\mathbb{N}}
\newcommand{\Z}{\mathbb{Z}}
\newcommand{\Q}{\mathbb{Q}}
\newcommand{\R}{\mathbb{R}}
\newcommand{\C}{\mathbb{C}}
\newcommand{\F}{\mathbb{F}}
\newcommand{\A}{\mathbb{A}}
\newcommand{\PP}{\mathbb{P}}
\newcommand{\E}{\mathbb{E}}
\newcommand{\eps}{\varepsilon}
\newcommand{\dd}{\,\mathrm{d}}
\begin{document}
\section*{UnsolvedMath ID: 20002789}
\textbf{Source problem code:} catalogue entry
\subsection{Definitions and mathematical statement}
The record contains three distinct source questions: develop relative endoscopy with the relative trace formula; interpret transfer from \(\mathrm{GL}_2\) to an elliptic torus, opposite the direction in MR3117742; and relate intertwining operators to character relations, with Labesse--Langlands as a model. The trailing ``Beyond Endoscopy'' is a following heading. No local field, relative spaces, transfer normalization, or packet hypotheses are specified.

The aid's directionality obstruction fixes a local field \(F\), a quadratic field extension \(E/F\), \(T=\operatorname{Res}_{E/F}\mathbb G_m\), and \(G=\mathrm{GL}_2\). Their connected dual groups are \(\widehat G=\mathrm{GL}_2(\C)\), \(\widehat T=(\C^\times)^2\), with Weil action on \(\widehat T\) swapping coordinates outside \(W_E\). For an L-homomorphism \({}^LG\to{}^LT\) respecting the Weil projection and algebraic on connected dual groups, the connected restriction must be
\[
 g\longmapsto((\det g)^m,(\det g)^m),\qquad m\in\Z.
\]
Characters of \(\mathrm{GL}_2\) factor through determinant, and Weil equivariance forces the two exponents equal. Therefore it cannot retract a standard torus-to-group L-embedding, whose connected restriction is \((z_1,z_2)\mapsto\operatorname{diag}(z_1,z_2)\), or distinguish torus parameters with equal product. This excludes that reverse L-homomorphism interpretation, not every opposite test-function operator.

A different conditional construction is transpose transfer. After choosing test-function/distribution spaces and Haar or spectral pairings, a forward kernel operator \(A:\mathcal H_T\to\mathcal H_G\) can have a transpose characterized by
\[
 \langle Af,h\rangle_G=\langle f,A^{\mathsf t}h\rangle_T.
\]
If represented by an integrable kernel \(K(g,t)\), then \((Af)(g)=\int K(g,t)f(t)\,\dd t\) and \((A^{\mathsf t}h)(t)=\int K(g,t)h(g)\,\dd g\), for bilinear pairings. A Hermitian adjoint instead conjugates the kernel. Such a reverse arrow depends on measures, domains, convergence, and the normalization of the forward stable-character kernel. The record supplies none of these analytic data, so it does not establish an unconditional reverse-transfer theorem or functorial inverse.

For the character relation, assume a finite abelian group \(A_0\) acts by an honest normalized family of self-intertwiners \(R(a)\), with \(R(a)R(b)=R(ab)\), commuting with a representation \(\pi\). Where all test-operator traces exist, define
\[
 D_a(f)=\operatorname{Tr}(\pi(f)R(a)),\qquad
 e_\chi=|A_0|^{-1}\sum_a\chi(a)^{-1}R(a),\qquad
 \Theta_\chi(f)=\operatorname{Tr}(\pi(f)e_\chi).
\]
Finite Fourier inversion gives
\[
 D_a=\sum_{\chi\in\widehat A_0}\chi(a)\Theta_\chi,\qquad
 \Theta_\chi=|A_0|^{-1}\sum_a\chi(a)^{-1}D_a.
\]
For \(|A_0|=2\), this is \(D_e=\Theta_++\Theta_-\), \(D_a=\Theta_+-\Theta_-\), the signed rank-one pattern invoked by the aid. Calling the \(\Theta_\chi\) packet characters requires the relevant constituent and multiplicity identification. A projective intertwiner action needs cocycle modifications and is not covered by ordinary Fourier inversion. These conditional formulas do not construct general relative endoscopy or identify all Labesse--Langlands transfer normalizations.
\subsection{Short English statement}
The aid separates three source questions. A reverse L-map from \(\mathrm{GL}_2\) to the elliptic torus collapses connected parameters to determinant data, while a transpose kernel transfer is a different conditional notion. Honest finite abelian intertwiner actions yield character relations by Fourier inversion, with packet normalization still requiring input.
\subsection{Sources}
\begin{itemize}
\item Local registry, permanent UnsolvedMath ID \(20002789\), exact public code \texttt{catalogue entry}; canonical location \texttt{aim-representation-theory-notes.json}, \texttt{notes[60]}. The AIM domain is Representation theory; the registry category is Algebra.
\item AIM, \emph{Automorphic Kernel Functions}, composite items (2), (3), (4); supplied reference \url{https://aimath.org/pastworkshops/automorphkernelproblems.pdf}. The following ``Beyond Endoscopy'' heading is not a fourth question. No PDF was accessed here. No proposer, year, remarks, or literature field is supplied.
\item MR3117742 for the torus-to-\(\mathrm{GL}_2\) transfer-factor direction, and Labesse--Langlands for intertwiner/character relations, as cited in the source. No titles, dates, theorem numbers, or full bibliographic entries are supplied.
\item Supplied research attempt 1: \texttt{partial\_result}, run and registry status \texttt{partially\_solved}, rigor label \texttt{proved}, low novelty confidence; classification \texttt{PARTIAL-PROGRESS}, statement status \texttt{reconstructed\_unverified}, created and updated 14 August 2026. The aid is machine generated, not peer reviewed, and requires independent verification before citation.
\end{itemize}
\end{document}
